\documentclass[12pt,a4paper]{article}
\usepackage[T2A]{fontenc}
\usepackage[utf8]{inputenc}
\usepackage[russian]{babel}
\usepackage{amsmath,amssymb}
\usepackage{array,booktabs,multirow,longtable,tabularx}
\usepackage[left=2cm,right=2cm,top=2cm,bottom=2cm]{geometry}
\usepackage{graphicx}
\usepackage{float}
\usepackage{caption}
\usepackage{tikz}
\usepackage{pgfplots}
\usepackage{hyperref}
\pgfplotsset{compat=1.18}
\usepackage{indentfirst}
\renewcommand{\arraystretch}{1.18}
\setlength{\tabcolsep}{6pt}
\hypersetup{colorlinks=true,urlcolor=blue,linkcolor=black}

\begin{document}

\noindent
\begin{minipage}[t]{0.6\textwidth}
    \vspace{0pt}
    Университет ИТМО \\
    Факультет \underline{\hspace{0.6cm} ФБИТ \hspace{5.95cm}}
\end{minipage}%
\hfill%
\begin{minipage}[t]{0.3\textwidth}
    \vspace{-14pt}
    \centering
    \begin{tikzpicture}
        \draw[black, line width=0.8pt] (0,0) rectangle (3.7,1.35);
        \node[font=\bfseries\Large] at (1.85,0.83) {ИТМО};
        \node[font=\scriptsize] at (1.85,0.35) {UNIVERSITY};
    \end{tikzpicture}
\end{minipage}
\rule{\textwidth}{1.5pt}
\vspace{0.01cm}

\noindent
Группа \underline{\hspace{1cm} 1.1 \hspace{2.85cm}} \hfill Работа проверена \underline{\hspace{4cm}}\\
Студент \underline{\hspace{0.2cm} Родыгин Даниил \hspace{0.2cm}} \hfill Работа выполнена \underline{\hspace{3.4cm}}\\
Преподаватель \underline{\hspace{0.935cm} Замолоцкий С.С. \hspace{1.5cm}} \hfill Отчет принят \underline{\hspace{4.84cm}}\\

\vspace{0.01cm}

\begin{center}
    \Large\textbf{Рабочий протокол и отчет} \\
    \Large\textbf{по проектной работе по физике}
\end{center}
\noindent\rule{\textwidth}{1.5pt}%
\vspace{-0.3cm}
\begin{center}
Интерактивная модель «Гравитационный катапульт»
\end{center}
\vspace{-0.3cm}
\noindent\rule{\textwidth}{1.5pt}

\section{Цель работы}
Разработка интерактивной компьютерной модели, демонстрирующей, как начальная скорость и угол запуска тела у поверхности планеты определяют форму его траектории: падение на планету, замкнутую орбиту, параболический уход или разомкнутую гиперболическую траекторию.

\section{Задачи}
\begin{enumerate}
    \item Реализовать математическую модель движения тела в центральном гравитационном поле планеты.
    \item Обеспечить возможность задания массы и радиуса планеты пользователем.
    \item Реализовать настройку начальной скорости и угла запуска с помощью слайдеров и прямого ввода чисел.
    \item Рассчитать первую и вторую космические скорости для выбранной планеты.
    \item Визуализировать планету, точку старта, спутник, прогнозную траекторию и фактический след движения.
    \item Отображать текущую скорость, высоту, ускорение свободного падения и удельную орбитальную энергию.
    \item Реализовать таблицу измерений, показывающую изменение физических величин во время полета.
    \item Добавить удобные элементы управления: запуск, паузу, перезапуск, сброс к параметрам Земли и переключатели слоев визуализации.
\end{enumerate}

\section{Объект исследования}
Объектом исследования является движение малого тела вблизи массивной планеты под действием силы всемирного тяготения. Планета считается гладким шаром радиуса \(R\) и массы \(M\), а тело рассматривается как материальная точка, масса которой не влияет на движение планеты.

\section{Метод исследования}
В работе использован метод математического и компьютерного моделирования. Движение тела рассчитывается в двумерной плоскости в центральном гравитационном поле. Пользователь задает параметры планеты и начальные условия запуска, после чего программа численно интегрирует уравнения движения и выводит результат в виде анимированной визуализации.

Для повышения устойчивости численного расчета используется метод velocity Verlet. Он лучше подходит для орбитальных задач, чем простой метод Эйлера, так как точнее сохраняет энергию системы на длительных промежутках времени.

\section{Рабочие формулы и исходные соотношения}
Гравитационный параметр планеты:
\begin{equation}
    \mu = GM,
\end{equation}
где \(G\) -- гравитационная постоянная, \(M\) -- масса планеты.

Расстояние от центра планеты до тела:
\begin{equation}
    r = \sqrt{x^2+y^2}.
\end{equation}

Ускорение тела в центральном гравитационном поле:
\begin{equation}
    \vec a = -\frac{GM}{r^3}\vec r.
\end{equation}

Модуль гравитационного ускорения:
\begin{equation}
    a = \frac{GM}{r^2}.
\end{equation}

Первая космическая скорость у поверхности планеты:
\begin{equation}
    v_1 = \sqrt{\frac{GM}{R}}.
\end{equation}

Вторая космическая скорость:
\begin{equation}
    v_2 = \sqrt{\frac{2GM}{R}}.
\end{equation}

Удельная орбитальная энергия:
\begin{equation}
    \varepsilon = \frac{v^2}{2} - \frac{GM}{r}.
\end{equation}

Начальные условия задаются из точки на поверхности планеты. При угле запуска \(\alpha\) к горизонту компоненты начальной скорости вычисляются как:
\begin{equation}
    v_x = v_0 \sin \alpha,
\end{equation}
\begin{equation}
    v_y = v_0 \cos \alpha.
\end{equation}

Критерий классификации траектории по энергии:
\begin{equation}
\begin{cases}
\varepsilon < 0, & \text{замкнутая эллиптическая орбита},\\
\varepsilon \approx 0, & \text{параболическая граница ухода},\\
\varepsilon > 0, & \text{разомкнутая гиперболическая траектория}.
\end{cases}
\end{equation}

\section{Описание разработанной модели}
Разработанный сайт представляет собой интерактивный симулятор «Гравитационный катапульт». На экране отображается планета, точка старта на ее поверхности, спутник и траектория движения. Пользователь может изменять параметры модели и сразу наблюдать, как меняется форма траектории.

В модели предусмотрены три характерных режима запуска:
\begin{enumerate}
    \item \(v_0 = v_1\) -- тело выходит на почти круговую орбиту.
    \item \(v_0 = v_2\) -- тело находится на границе ухода из гравитационного поля планеты.
    \item \(v_0 > v_2\) -- тело движется по разомкнутой траектории.
\end{enumerate}

Если выбранная скорость слишком мала или угол направляет тело к поверхности, модель показывает пересечение траектории с планетой и классифицирует результат как столкновение.

\section{Функционал интерфейса}
В пользовательском интерфейсе реализованы следующие элементы:
\begin{enumerate}
    \item \textbf{Слайдер начальной скорости} с возможностью прямого ввода значения в км/с.
    \item \textbf{Слайдер угла запуска} с возможностью прямого ввода угла в градусах.
    \item \textbf{Поля массы и радиуса планеты}, позволяющие моделировать не только Землю, но и другие планеты.
    \item \textbf{Кнопки предустановленных скоростей}: \(v_1\), \(v_2\), \(1.25v_2\).
    \item \textbf{Кнопки управления симуляцией}: запуск, пауза, перезапуск и сброс к параметрам Земли.
    \item \textbf{Переключатели визуальных слоев}: след траектории, вектор скорости, вектор гравитации и автомасштаб.
    \item \textbf{Информационные карточки}, показывающие скорость, высоту и орбитальную энергию.
    \item \textbf{Таблица измерений}, фиксирующая изменение времени, высоты, скорости, энергии и ускорения во время полета.
\end{enumerate}

\section{Использованные технологии}
\begin{longtable}{p{0.27\textwidth}p{0.65\textwidth}}
\toprule
\textbf{Технология} & \textbf{Назначение в проекте}\\
\midrule
React & построение интерфейса сайта и управление состоянием симуляции.\\
TypeScript & строгая типизация физических параметров, телеметрии и настроек модели.\\
Vite & быстрый запуск локального сервера и сборка фронтенд-проекта.\\
Canvas API & отрисовка планеты, спутника, звездного фона, траектории и векторов.\\
p2-es & физический слой для представления динамического тела и мира без стандартной гравитации.\\
GSAP & анимации появления интерфейса и световой эффект при запуске.\\
lucide-react & набор иконок для кнопок, карточек и панелей управления.\\
CSS & адаптивная верстка, оформление слайдеров, переключателей, таблицы и визуальных панелей.\\
\bottomrule
\end{longtable}

\section{Педагогическая ценность модели}
Разработанный продукт может использоваться как учебный инструмент при изучении темы движения тел в гравитационном поле. Его педагогическая ценность состоит в следующем:
\begin{enumerate}
    \item Пользователь может \textbf{наглядно увидеть} различие между круговой, эллиптической, параболической и гиперболической траекторией.
    \item Формулы первой и второй космической скорости связываются с \textbf{визуальным результатом} на экране.
    \item Изменение скорости и угла запуска позволяет \textbf{экспериментально проверить}, при каких условиях тело остается на орбите, падает или покидает планету.
    \item Таблица измерений помогает связать анимацию с конкретными физическими величинами: скоростью, высотой, ускорением и энергией.
    \item Модель можно расширять: добавлять атмосферное сопротивление, несколько планет, масштабирование массы спутника или сравнение разных численных методов.
\end{enumerate}

\section{Окончательные результаты}
В результате выполнения проектной работы создан интерактивный сайт, который моделирует:
\[
r(t),\quad v(t),\quad a(t),\quad \varepsilon(t),\quad v_1,\quad v_2
\]
и демонстрирует зависимость формы траектории от начальной скорости и угла запуска.

Основной содержательный результат работы состоит в подтверждении следующего вывода:
\[
\text{форма траектории тела определяется его начальной энергией и направлением скорости.}
\]

При скорости, близкой к первой космической, тело может двигаться по замкнутой орбите. При скорости, близкой ко второй космической, тело находится на границе ухода. При скорости больше второй космической траектория становится разомкнутой.

\section{Выводы и анализ результатов работы}
В ходе выполнения проектной работы была разработана и проанализирована интерактивная модель движения тела в гравитационном поле планеты. В модели реализованы все основные требования: визуализация планеты и спутника, запуск с поверхности, настройка начальной скорости и угла, расчет первой и второй космической скорости, вывод текущей орбитальной энергии и таблица изменения физических величин.

Физическая часть проекта основана на законе всемирного тяготения и понятии удельной орбитальной энергии. Проведенное моделирование показывает, что одна и та же планета может давать принципиально разные типы движения в зависимости от начальной скорости: при недостаточной скорости тело падает, при скорости порядка первой космической удерживается на орбите, а при достижении второй космической получает возможность уйти от планеты.

Практическая значимость работы состоит в том, что полученный сайт позволяет изучать тему не только по формулам, но и через непосредственное взаимодействие с моделью. Пользователь может менять параметры, запускать сценарии и сразу видеть результат на анимированной сцене и в таблице. Благодаря этому проект выступает как полноценное цифровое учебное пособие по теме космических скоростей и орбитального движения.

\end{document}
